\documentclass[11pt,a4paper]{article}

\usepackage[margin=1in]{geometry}
\usepackage[T1]{fontenc}
\usepackage{listings}
\usepackage{xcolor}
\usepackage{booktabs}
\usepackage{enumitem}
\usepackage[hidelinks]{hyperref}

\definecolor{codebg}{gray}{0.96}
\definecolor{codekw}{RGB}{0,90,160}
\definecolor{codecm}{RGB}{100,110,100}

% Defined explicitly rather than relying on listings' built-in Rust support,
% which is absent from older distributions and errors out when requested.
\lstdefinelanguage{Rust}{
  keywords={as,async,await,break,const,continue,crate,dyn,else,enum,extern,
    false,fn,for,if,impl,in,let,loop,match,mod,move,mut,pub,ref,return,self,
    Self,static,struct,super,trait,true,type,unsafe,use,where,while},
  sensitive=true,
  morecomment=[l]{//},
  morecomment=[s]{/*}{*/},
  morestring=[b]",
}

\lstdefinestyle{rust}{
  language=Rust,
  backgroundcolor=\color{codebg},
  basicstyle=\ttfamily\small,
  keywordstyle=\color{codekw}\bfseries,
  commentstyle=\color{codecm}\itshape,
  breaklines=true,
  frame=none,
  xleftmargin=1em,
  framexleftmargin=1em,
  showstringspaces=false,
  captionpos=b
}
\lstset{style=rust}

% Deliberately takes no argument: the braces at each call site become an
% ordinary group, which keeps \verb usable inside the note text.
\newcommand{\note}{\par\smallskip\noindent\textit{Why it matters:} }

\title{Horizon Function Map\\[0.3em]\large Deferred Scope and Known Exclusions}
\author{}
\date{26 July 2026}

\begin{document}
\maketitle

\section{Purpose of this document}

This is a record of everything discussed and \emph{deliberately set aside} while
designing the function map. It exists so that a decision made for good reasons
today is not mistaken for an oversight later, and so that the cost of each
exclusion is written down while it is still fresh.

Nothing here is a bug. Each item is either an accepted limitation, a deferral, or
an alternative that was evaluated and passed over. Items that remain genuinely
undecided are collected separately in Section~\ref{sec:open}.

Where a number appears, it was measured on real code --- Horizon's own
\verb|src/|, the purpose-built demo crates, or published benchmarks --- rather
than estimated.

\textbf{Reading note (post-implementation).} Several design-phase assumptions
were later reversed (no compile gate; crates as spine nodes; allowlisted macro
recovery; correctness harness). This file has been reconciled with the
specification of record in
\texttt{docs/requirements/rust-function-map.md}. Prefer that document when the
two disagree on \emph{current} behaviour; this file keeps the measured costs and
the reasoning that led to each exclusion.

\section{The frame these exclusions follow from}

Four decisions shape everything below.

\begin{description}[style=nextline,leftmargin=1em]
  \item[Functions only.] The map contains free functions, their calls, and
    documentation. Data structures are not map nodes.
  \item[Free functions only.] Methods and \verb|impl| / \verb|trait| items are
    \textbf{permanently out of scope} for this tool (declarations and call
    sites). Associated functions on types are excluded for the same reason.
  \item[Build the tree always.] There is no compile gate. Mid-edit and
    non-compiling source still yields a map; conflicts and unresolved markers
    make incompleteness visible. (An earlier ``valid code only'' frame was
    rejected --- see Section~\ref{sec:gates}.)
  \item[Batch updates.] The map is a single-shot batch analysis. Live
    incremental updating while typing is a non-goal. The host product's
    rebuild trigger (when a codebase is ``uploaded'', whether prior JSON is
    retained) is unspecified; the tool itself has no cache.
\end{description}

\section{Deferred by explicit decision}

\subsection{Methods and \texttt{impl} blocks}

Excluded entirely: method declarations inside \verb|impl| blocks, and calls of
the form \verb|x.foo()|. Permanently out of scope --- not a ``phase~2'' of this
resolver.

\note{This is the largest single exclusion by volume. Of Horizon's \textbf{1{,}394
call sites, 1{,}053 are method calls} --- roughly 75\%. Only 266 are the
path-form calls the map currently describes. The map will therefore represent a
minority of the call activity in a typical Rust codebase.}

Resolving them requires knowing the type of the receiver, which needs type
inference rather than name lookup. Two sub-cases are known to be unresolvable
without a real type checker and would remain excluded even in a separate
type-aware product:

\begin{itemize}[nosep]
  \item receivers typed through a generic parameter (\verb|T: Store|),
  \item receivers behind a trait object or opaque type (\verb|&dyn Store|,
    \verb|impl Store|).
\end{itemize}

A related correctness rule was identified for any future method work: Rust probes
inherent methods before trait methods, so when a type has both, the inherent one
wins. The demo fixture \verb|tests/fixtures/method-ambiguity/| contains exactly
this collision.

\begin{lstlisting}
impl Cache {
    pub fn get(&self, key: &str) -> Option<&Item> { ... }
}

impl Store for Cache {
    fn get(&self, key: &str) -> Option<&Item> { ... }
}
\end{lstlisting}

\subsection{Associated functions and constructors}

Also permanently out of scope, counted in summary counters rather than drawn:

\begin{itemize}[nosep]
  \item associated functions on types (\verb|Vec::new|, \verb|Type::assoc|) ---
    \verb|associated_dropped|;
  \item enum variants and tuple-struct constructors (\verb|Ok(x)|,
    \verb|CallTarget::Resolved(...)|) --- \verb|constructor_dropped|.
\end{itemize}

\note{Both are syntactically \verb|CallExpr| in \verb|ra_ap_syntax|. Recording
them as \verb|Unresolved| would mis-diagnose a type or constructor as a missing
free function. Classification uses collected local type names, prelude variants,
primitives, and a strict UpperCamelCase convention heuristic when facts are
absent (failures surface as \verb|Unresolved|, not silent drops).}

\subsection{Data structures}

Excluded from the \emph{emitted} map: \verb|struct|, \verb|enum|, and
\verb|trait| definitions as containment nodes; the use of one data structure
inside another's definition; and field types as map edges.

\note{Type \emph{names} (and enum variants) are collected internally so
constructors and associated functions can be classified and dropped. That index
is not part of the JSON map. Representing data structures as first-class map
nodes remains deferred, not abandoned.}

\subsection{Compile gate (rejected)}

An earlier frame required the map to update only from code that compiles. That
was replaced: build always, mark conflicts / unresolved.

\note{Error tolerance was one of the original reasons for choosing
\verb|ra_ap_syntax| over \verb|syn|. Keeping a compile gate would have thrown
away that capability and blocked mid-edit repositories. Ambiguous globs
(\verb|E0659|) and missing names (\verb|E0425|) are first-class map outcomes
under the never-guess rule.}

Validity-gate mechanisms that were evaluated and \emph{not} adopted are recorded
in Section~\ref{sec:gates} for history.

\subsection{Live incremental updating}

Excluded: updating the map continuously as code is typed.

\note{The tool is a single-shot batch analyser with no incremental state. The
original ``update while typing'' requirement conflicted with batch rebuilds and
is a non-goal. Host-product storage and rebuild policy remain open
(requirements, Open questions).}

\subsection{Call categories excluded from the map}

Kinds of call recognised by the analyser but absent from the tree (except where
noted).

\begin{description}[style=nextline,leftmargin=1em]

  \item[External calls.] Calls into \verb|std| / \verb|core| / \verb|alloc| /
    \verb|proc_macro|, or into a Cargo dependency whose kind is external.
    Deliberately dropped --- not drawn as \verb|Unresolved|, not given a new
    variant. Counted in \verb|MapSummary.external_dropped|.
    \note{Someone reading the never-silently-missing rule literally might try to
    ``restore'' these as unresolved sites. Do not: external callees are outside
    the indexed universe by design. A function whose body only calls
    \verb|std| appears as a childless leaf; that is intentional.}

  \item[Indirect calls.] A call through a variable holding a function.
\begin{lstlisting}
fn apply(s: &str, f: fn(&str) -> String) -> String {
    f(s)                     // `f` is a parameter, not a function name
}
let k = |x: &str| normalize(x);
k(label)                     // `k` is a local closure binding
\end{lstlisting}
    Which function runs depends on what the caller passed and may differ per
    invocation. Determining it requires dataflow analysis.

  \item[Functions passed as values.] In \verb|apply(&n, normalize)| the name
    \verb|normalize| is mentioned but not called.
    \note{This relationship vanishes completely today. A function can be passed
    throughout a codebase and still appear unreferenced in the map.}

\end{description}

\paragraph{Macro-hidden calls (partially recovered).}
 The parser leaves macro
arguments as unstructured token trees, so \verb|mean(v)| inside
\verb|format!("{}", mean(v))| is invisible to a normal tree walk. Measured on
Horizon this hid on the order of tens of path-form calls against a few hundred
visible ones.

\textbf{Implemented:} an allowlist of std/core expression macros
(\verb|format!|, \verb|println!|, \verb|assert_eq!|, \verb|vec!|, \ldots) is
opened by re-parsing the token-tree interior; recovered sites resolve through
the normal path and set \verb|CallSite.from_macro|. Nested allowlisted macros
are opened recursively (depth-bounded).

\textbf{Still deferred / deliberately opaque:} user-defined macros,
\verb|matches!| / pattern-position macros, \verb|stringify!| / \verb|concat!|,
\verb|cfg!|, \verb|quote!|, \verb|macro_rules!| / \verb|macro| definition
bodies, attribute / derive token trees. Prefer a miss over a fabricated edge.

\paragraph{Macro-hidden \texttt{mod} declarations (allowlisted item reparse).}
 Allowlisted item-pasting macros (\verb|cfg_if!|, Tokio-style names starting
\verb|cfg_|) hide \verb|mod| declarations inside token trees. Recovery
re-parses the \emph{invocation} token tree (depth-bounded at 8) and walks the
resulting modules. All \verb|cfg_if!| branches are \textbf{unioned} --- the tool
does not choose a build configuration, for the same reason it keeps
\verb|#[cfg]|-duplicated functions.

\textbf{Not recovered:} serde's \verb|crate_root!()| pattern. The modules live
in a \verb|macro_rules!| \emph{definition} body; the call site is empty.
Expanding definition bodies is a different and riskier problem than re-parsing
an invocation's tokens. Prefer a missing subtree over a fabricated module tree.

\textbf{First-wins on cfg-twinned \texttt{\#[path]}:} when two \verb|#[cfg]| arms
declare the same module name with different \verb|#[path]| targets, the first
declaration wins --- one module path maps to one file. Real limitation.

\section{Language features knowingly unhandled}

\subsection{Configuration-gated duplicate definitions}

\begin{lstlisting}
#[cfg(unix)]    fn open() { }
#[cfg(windows)] fn open() { }
\end{lstlisting}

One path, two real definitions. Which is live depends on a build configuration
nobody has chosen. Horizon emits a \verb|Conflict| naming every candidate
(FunctionIds get \verb|#L{line}| suffixes). That is the never-guess outcome,
not a silent pick.

\subsection{Modules absent from the tree spine}

The map's containment hierarchy is crate $\rightarrow$ folder $\rightarrow$
file $\rightarrow$ function. Modules are not nodes, on the working assumption
that a module corresponds to a file. That assumption was tested against the
compiler and is false in four distinct ways:

\begin{table}[h]
\centering
\begin{tabular}{@{}lll@{}}
\toprule
\textbf{Case} & \textbf{Module} & \textbf{File} \\
\midrule
Ordinary            & \verb|text|              & \verb|text.rs| \\
Inline module       & \verb|text::case|        & none --- written inside \verb|text.rs| \\
Undeclared file     & none                     & \verb|orphan.rs| --- never compiled \\
\verb|#[path]|      & \verb|tidy_name|         & \verb|renamed_on_disk.rs| \\
Crate root          & the root itself          & \verb|lib.rs| --- contributes no segment \\
\bottomrule
\end{tabular}
\end{table}

\note{Two consequences follow. Inline modules disappear from the spine, so a file
may contain two functions of the same name at different module depths --- and a
call node that references only its declaring \emph{file} cannot distinguish them
(resolved targets therefore reference \emph{function} ids). Second, a file
declared nowhere is not part of the crate at all. File discovery is
declaration-driven (\verb|mod| walk), never a directory glob.}

A further consequence, verified: a \verb|mod| inside a file nests rather than
merging, so a file \verb|selfdecl.rs| declaring \verb|pub mod selfdecl| yields the
path \verb|crate::selfdecl::selfdecl::inner_fn|.

\subsection{Textual inclusion}

\verb|include!("other.rs")| splices a file's text in with no module boundary.
Rare outside code generation, but it means even a complete \verb|mod| walk can
miss source that is genuinely compiled.

\subsection{Scope-local \texttt{use} declarations}

\verb|use| is legal inside functions and modules, not only at file top level.
A \verb|use| inside a function body is attributed to the enclosing
\emph{module} (module-wide) --- wider than rustc. Module-level \verb|use|
(including inside inline \verb|mod| blocks) is attributed to that module
correctly. The widening is a documented limitation; narrowing body-scoped
imports would change resolution outcomes and is deferred.

\subsection{Glob imports into dependencies}

\verb|use some_dependency::*;| is not specially expanded across crates.
Within-crate globs are resolved against the index with visibility-filtered
candidate sets. Cross-crate edges use explicit paths / imports into declared
path dependencies.

\subsection{Traits imported as \texttt{\_}}

\verb|use SomeTrait as _;| brings a trait into scope for method resolution
without binding a nameable identifier, so the trait becomes invisible to the
index. Method-related, and therefore doubly out of scope.

\subsection{Visibility asymmetry (intentional)}

Three distinct policies --- do not ``unify'' them:

\begin{itemize}[nosep]
  \item \textbf{Direct within-crate edges:} visibility is \emph{not} filtered.
    A path the author wrote still appears even if rustc would emit
    \verb|E0603|.
  \item \textbf{Glob candidate sets:} visibility \emph{is} filtered so private
    items do not manufacture \verb|E0659| conflicts rustc would never report.
  \item \textbf{Cross-crate reachability:} visibility \emph{is} enforced
    (\verb|pub| item behind an all-\verb|pub| module chain) plus the
    dependency gate.
\end{itemize}

\section{Build-system cases knowingly unhandled}

Crate discovery is driven by \verb|cargo metadata --no-deps --offline|, which
was measured at \textbf{322\,ms on Horizon} and performs no compilation. Path
dependencies outside the workspace are found by recursing into each
\verb|path| directory. The following remain unhandled.

\begin{description}[style=nextline,leftmargin=1em]
  \item[\texttt{[patch]} and \texttt{[replace]}.] These can redirect an apparent
    registry dependency to a local directory, so a crate can be local without its
    own entry saying \verb|path|. Same for vendored setups where
    \verb|.cargo/config.toml| places every dependency on disk.
  \item[Optional and feature-gated dependencies.] Presence depends on the
    selected feature set, which is genuinely undecidable without choosing one.
  \item[Target-specific dependencies.] Entries under
    \verb|[target.'cfg(unix)'.dependencies]| vary by platform.
  \item[Development dependencies.] Visible only to tests, examples, and benches,
    not to the library --- filtered out of the dependency list used for
    resolution. Scale validation shows the cost: test-only externals surface as
    \verb|Unresolved| rather than \verb|external_dropped|. Including them as
    external roots (not walked) remains an open product decision.
  \item[Examples, integration tests, benches, \texttt{build.rs}.]
    Deliberately \emph{not} discovered as map crates. They are secondary
    surfaces; including them would inflate maps and hit the same
    \verb|[dev-dependencies]| blind spot. Not an accidental omission of target
    kinds.
  \item[Non-Cargo projects.] Builds driven by \verb|rust-project.json|, Buck, or
    Bazel are out of scope; everything assumes Cargo.
\end{description}

\paragraph{Handled (recorded because once surprising).} Package names and
editions come from \verb|cargo metadata|, not directory heuristics. A binary in
\verb|src/bin/| is a separate crate node from the library beside it and reaches
it by package name (implicit path dependency), exactly as Cargo does.
FunctionIds disambiguate same-named lib+bin pairs as \verb|{name}[bin]|.
Manifest dependency renaming
(\verb|alias = { package = "real-name", path = "..." }|) is supported: the
import name is what code uses; the package name identifies the path dep.
Library-like target kinds include \verb|proc-macro| (as well as \verb|lib| /
\verb|rlib| / \verb|dylib| / \verb|cdylib| / \verb|staticlib|) ---
\verb|cargo metadata| reports proc-macro crates with kind \verb|proc-macro|,
not \verb|lib|, so omitting that kind silently dropped crates such as
\verb|serde_derive|.

\section{Alternatives evaluated and set aside}

\subsection{rust-analyzer SCIP indexing}
\label{sec:scip}

\verb|rust-analyzer scip .| emits a complete semantic index --- definitions,
references, resolved cross-references, hover documentation --- using real name
resolution and type inference rather than heuristics. Published measurements:

\begin{table}[h]
\centering
\begin{tabular}{@{}lrrl@{}}
\toprule
\textbf{Codebase} & \textbf{Time} & \textbf{Peak memory} & \textbf{Result} \\
\midrule
ruff              & 73\,s   & 2.1\,GB & 646k occurrences, 81\% resolved \\
ripgrep           & 35\,s   & 1.1\,GB & full index \\
rust-analyzer     & 35\,s   & ---     & full index \\
mozilla-central   & 20\,min & ---     & full index \\
\midrule
ruff, via live LSP & 254\,s & 15.4\,GB & 9\% resolved, degraded \\
\bottomrule
\end{tabular}
\end{table}

Unresolved references are \verb|std| and dependencies, marked
\verb|EXTERNAL_SYMBOL| --- which is the intended external-dependency rule,
implemented for free. Coverage includes methods.

\note{This resolves roughly 81\% including methods, against the hand-built
approach's roughly 55\% of path-form calls with methods excluded. The decision
is to build our own for now regardless. The costs of adopting SCIP are a required
toolchain (\verb|rustup component add rust-analyzer|), a batch snapshot rather
than incremental output, and (at decision time) unverified handling of comments
as first-class nodes. LSIF from rust-analyzer is used only as a
\emph{correctness oracle} for free-function edges on Horizon itself --- not as
the production indexer.}

\subsection{Validity gates not chosen}
\label{sec:gates}

Three mechanisms for establishing that code compiles were identified. Under the
``build always'' decision they are unused; recorded for history.

\begin{itemize}[nosep]
  \item \verb|cargo check --message-format=json| --- the authority. One JSON
    object per line; the \verb|cargo_metadata| crate parses the stream. Stops
    before code generation, so a small class of codegen-time diagnostics never
    appears.
  \item \verb|rust-analyzer diagnostics <path>| --- purpose-built as a gate;
    exits non-zero if any error is found.
  \item \verb|rust-analyzer unresolved-references <path>| --- narrower and
    aimed exactly at this problem, reporting \verb|file:line:column: text|. Known
    false positives around procedural macros unless the workspace is loaded with
    \verb|all_targets: true|.
\end{itemize}

\paragraph{Rejected outright: parser errors as a validity gate.} Syntactic
validity and compilability are almost unrelated. Of the compiler errors produced
while testing these rules, only one was syntactic --- a dash in a path. Every
other one parsed perfectly and failed later:

\begin{lstlisting}[language={}]
E0433  cannot find module `text`        <- parsed fine
E0425  cannot find function `get`       <- parsed fine
E0603  module `secret` is private       <- parsed fine
E0659  `get` is ambiguous               <- parsed fine
E0308  mismatched types                 <- parsed fine
\end{lstlisting}

\section{Known weaknesses in what already exists}

\begin{description}[style=nextline,leftmargin=1em]
  \item[Correctness outside the measured set.] A standing harness reports zero
    false positives on fixture oracles and on an LSIF comparison against
    Horizon itself (see \texttt{docs/correctness-measurement.md}). That does
    \emph{not} speak for large or unfamiliar corpora at LSIF grade (a hand
    sample on ripgrep / Cellular Automata is recorded in
    \texttt{docs/scale-validation.md}).
  \item[UpperCamelCase drop heuristic.] Unknown path segments matching strict
    UpperCamelCase (ASCII uppercase start, alphanumeric only, at least one
    lowercase) may be classified as types and dropped. Segments that
    \emph{fail} that test surface as \verb|Unresolved| rather than a silent
    drop. Modules present in the module tree are unaffected; UpperCamelCase
    modules \emph{absent} from the tree can still be mis-dropped.
  \item[Nested free-function unqualified recursion.] An unqualified recursive
    call inside a nested \verb|fn| can surface as \verb|Unresolved|. No
    dedicated fixture yet.
  \item[Function-body \texttt{use} widening.] Documented; wider than rustc.
  \item[Re-export hop bound.] Chains are followed at most 32 hops
    (\verb|REEXPORT_HOP_LIMIT|), then \verb|Unresolved|.
  \item[Consumer-side cross-crate globs.] \verb|use dep::*| in the
    \emph{calling} crate is still not expanded. Foreign crates' own
    \verb|pub use dep::*| facades \emph{are} followed (settled; see
    requirements Cross-crate reachability).
\end{description}

\section{Undecided}
\label{sec:open}

These remain genuinely open (aligned with the requirements Open questions).
Decisions that were open at design time and are now settled are listed at the
end of this section for the record.

\begin{enumerate}[leftmargin=1.5em]
  \item \textbf{Update trigger.} Host-product policy for when to rebuild;
    the tool has no cache (see requirements).
  \item \textbf{Parallelism.} Pipeline is sequential; scale validation shows
    ~0.8--1.1\,MB/s analysed on dense trees and ~10--20\,min for a 50k-file
    estimate --- parallelism is the obvious lever.
  \item \textbf{Broader macro recovery.} User macros and definition-body
    expanders stay opaque; any opt-in needs its own false-positive measurement.
  \item \textbf{Cross-crate LSIF / large-corpus precision.} Unmeasured beyond
    purpose-built multi-crate fixtures and a hand sample on unfamiliar code.
  \item \textbf{Dev-dependencies as external roots.} Whether to list them so
    test-only calls become \verb|external_dropped| rather than
    \verb|Unresolved|.
\end{enumerate}

\paragraph{Settled since the design draft (no longer open).}
\begin{itemize}[nosep]
  \item Doc comments attach to functions (outer) and files (inner).
  \item Crates are spine nodes above folders.
  \item Call edges reference function ids, not files.
  \item No validity gate; edition comes from \verb|cargo metadata|.
  \item File discovery is a \verb|mod| walk, not a directory glob.
  \item Manifest dependency renaming is supported.
  \item Proc-macro (and other library-like) target kinds are discovered;
    examples / tests / benches / \verb|build.rs| are deliberately excluded.
  \item Allowlisted item-macro \verb|mod| recovery (\verb|cfg_if!| / \verb|cfg_*|)
    with union-of-branches; \verb|crate_root!()|-style definition bodies stay closed.
  \item Capitalisation check is strict UpperCamelCase; failures surface as
    \verb|Unresolved| rather than silent drops.
  \item Correctness harness shares the extract/resolve pipeline with the CLI
    (no parallel twin walk).
  \item Cross-crate facade re-exports (\verb|pub use| / renamed / glob /
    multi-crate chains) are followed with the same hop bound as within-crate
    re-exports; initial entry remains dependency-gated.
\end{itemize}

\section{Original goals not addressed today}

\begin{itemize}[nosep]
  \item \textbf{Multi-language support.} The project is intended to work on
    arbitrary codebases; everything designed today is Rust-specific.
  \item \textbf{Methods / type-aware resolution.} Permanently out of this
    tool's scope; would be a separate product.
  \item \textbf{Identity under re-export.} One definition can be reachable by
    several paths. Import aliases and bounded re-export hops are handled;
    a general alias table as a first-class map feature is not.
\end{itemize}

\end{document}
